% MAT670 -- Interlude E: random packings, entropy and information.
% Entirely new material (not in the handwritten notes). Labels prefixed SM:.
% Includes the coordinator's section chapters/tetra.tex (labels TET:) verbatim.

\chapter[Interlude E: Entropy and information]{Interlude E: Random packings, entropy and information}
\label{SM:ch}

\begin{addedblock}
This interlude is new: it is not in the handwritten notes. It picks up three threads that the course
left loose: the coin flips and the gas in a box of Lecture~7, the Edwards entropy and the Fibonacci
words of Lecture~11, and the project topics ``error correcting codes and information theory'' and
``entropy of jammed systems'' listed at the end of Lecture~6. To keep it readable it is set in black;
all of it should be read as added material.
\end{addedblock}

Most of this course was about the \emph{best} configurations: the densest packing, the largest
kissing configuration, the maximum of the injectivity radius. Lecture~7 already suggested a second
point of view. For large $N$ one does not ask for the best configuration but for the \emph{typical}
one, and one answers with Stirling's formula. Lecture~11 then did exactly this for jammed discs in a
channel, and found the golden ratio.

In this interlude we follow that second point of view: hard spheres as statistical mechanics with no
energy at all (Section~\ref{SM:sec:hard}); random and jammed packings (Section~\ref{SM:sec:random}); why
dense liquids do not simply crystallize (Sections~\ref{SM:sec:frustration}--\ref{SM:sec:curved}); Edwards'
statistical mechanics of the Lecture~11 channel, solved completely (Section~\ref{SM:sec:edwards}); and
information theory (Section~\ref{SM:sec:info}), where the channel's maximal Edwards entropy turns out to
be a Shannon capacity and $E_8$ and Leech come from perfect codes.

Throughout, $\ln$ is the natural logarithm and $\log_2$ the logarithm to base $2$. Boltzmann's constant
is set to $1$.

%======================================================================
\section{Hard spheres as statistical mechanics}
\label{SM:sec:hard}

Recall the remark at the end of \S\ref{L11:sec:boltzmann}: for hard particles the energy of a
configuration is $0$ if no two particles overlap and $+\infty$ if two do, so the Boltzmann weight
$e^{-\beta E}$ is the indicator function of the admissible configurations, at every temperature.
For $N$ balls of diameter $1$ with centres in a region $\Omega\subset\R^d$ of volume $V$, the
admissible configurations form the \emph{hard-sphere configuration space}
\[
  \Conf_{\ge1}(N,\Omega)=\{(x_1,\dots,x_N)\in\Omega^N:\ \abs{x_i-x_j}\ge1\ \text{for all }i\ne j\},
\]
a closed subset of the configuration space $\Conf(N,\Omega)$ of Lecture~7
(Definition~\ref{L07:def:conf}). On the sphere it is the space $\Conf(N,\theta)$ of
Definition~\ref{L10:def:conftheta}, whose topology we computed with min-type Morse theory. The
(configurational) \emph{partition function} is its volume,
\[
  Z_N(\Omega)=\frac1{N!}\vol\Conf_{\ge1}(N,\Omega),
\]
the factor $1/N!$ accounting for the indistinguishability of the particles. The \emph{entropy} is
$S=\ln Z_N$, and since the energy is identically zero the free energy is $F=-TS$. Temperature only
sets a scale. The \emph{pressure} is
\[
  \beta P=\frac{\partial S}{\partial V},
\]
the rate at which the configuration space grows when the box grows. (Change $\Omega$ by a homothety
to make this a derivative in one variable.) So for hard spheres \emph{everything is entropy}. This
is worth remembering in Section~\ref{SM:sec:edwards}, where Edwards runs the same dictionary in the
opposite direction and treats volume as if it were energy.

\subsection*{The Kepler problem is about where the free energy is finite}

Let $\dens_d$ be the maximal packing density of balls in $\R^d$, and let
\[
  \rho_{\mathrm{cp}}=\frac{\dens_d}{\vol\Ball^d(\tfrac12)}
\]
be the corresponding maximal \emph{number density} of centres for balls of diameter~$1$. In
dimension~$3$, Hales' theorem gives $\dens_3=\pi/\sqrt{18}$ and so $\rho_{\mathrm{cp}}=\sqrt2$, the
number density of the face-centred cubic lattice with nearest-neighbour distance~$1$.

\begin{proposition}\label{SM:prop:kepler}
Let $Q_\ell=[0,\ell]^d$, and let $N=N(\ell)$ with $N/\ell^d\to\rho$ as $\ell\to\infty$.
\begin{enumerate}[nosep]
\item If $N>\rho_{\mathrm{cp}}(\ell+1)^d$, then $Z_N(Q_\ell)=0$. In particular $Z_N(Q_\ell)=0$ for all
  large $\ell$ when $\rho>\rho_{\mathrm{cp}}$.
\item If $\rho<\rho_{\mathrm{cp}}$, then $\liminf_{\ell\to\infty}\frac1N\ln Z_N(Q_\ell)>-\infty$.
\end{enumerate}
\end{proposition}

\begin{proof}
(1) Given $N$ admissible centres in $Q_\ell$, the balls of radius $\frac12$ around them lie in the
cube $[-\frac12,\ell+\frac12]^d$ of side $\ell+1$. Repeat this cube periodically. The result is a
periodic packing of density $N\vol\Ball^d(\frac12)/(\ell+1)^d$, which is at most $\dens_d$.

(2) Choose $\rho<\rho''<\rho_{\mathrm{cp}}$. There is a periodic packing of balls of diameter $1$
with some number density $\rho'>\rho''$, since periodic packings come arbitrarily close to the optimal density
(Exercise~\ref{L1:ex:periodic-approx} of Lecture~1). Scale it by $\lambda=(\rho'/\rho'')^{1/d}>1$. The
centres now have number density $\rho''$ and are at mutual distance at least $\lambda$, so each centre may
be moved independently anywhere in a ball of radius $r=(\lambda-1)/2$ without creating an overlap.
The cube $Q_\ell$ contains $\rho''\ell^d-O(\ell^{d-1})$ of these wiggle balls, which is more than $N$ for
large $\ell$; keep $N$ of them. Assigning the $N$ labelled particles to the $N$ wiggle balls in any of
the $N!$ ways gives disjoint pieces of $\Conf_{\ge1}(N,Q_\ell)$, so
$Z_N\ge\frac1{N!}N!\,(\vol\Ball^d(r))^N$ and $\frac1N\ln Z_N\ge\ln\vol\Ball^d(r)$.
\end{proof}

The limit $s(\rho)=\lim\frac1N\ln Z_N$ exists for $\rho<\rho_{\mathrm{cp}}$ (subadditivity; see Ruelle's
book), so the free energy is finite at every density below $\rho_{\mathrm{cp}}$ and infinite at every
density above it. At $\rho_{\mathrm{cp}}$ itself one expects $s=-\infty$ and a divergent pressure, as for the
Tonks gas below, but the proposition does not decide this. Thus the theorems of Hales in $\R^3$, and of Viazovska and
Cohn--Kumar--Miller--Radchenko--Viazovska in $\R^8$ and $\R^{24}$, determine the density at which the
free energy of the hard-sphere gas becomes infinite.

\subsection*{One dimension: the Tonks gas}

On a line everything can be computed. Lecture~11 (\S\ref{L11:sec:hardrods}) found that the gaps
between $n$ rods of length $1$ in an interval of length $\ell$ form a simplex of volume
$(\ell-n)^n/n!$. Ordering the labels gives a factor $n!$, so
\[
  Z_n(\ell)=\frac{1}{n!}\cdot n!\cdot\frac{(\ell-n)^n}{n!}=\frac{(\ell-n)^n}{n!}.
\]

\begin{proposition}[Tonks 1936]\label{SM:prop:tonks}
The hard-rod gas has equation of state
\[
  \beta P=\frac{\partial}{\partial\ell}\ln Z_n(\ell)=\frac{n}{\ell-n}=\frac{\rho}{1-\rho},\qquad\rho=\frac n\ell,
\]
and entropy per particle $\frac1n\ln Z_n\to1+\ln\frac{1-\rho}{\rho}$ as $n,\ell\to\infty$ with
$n/\ell\to\rho\in(0,1)$.
\end{proposition}

The first formula is a one-line computation, and the second is Stirling's formula. Compare
Exercise~\ref{L11:ex:tonks}, where the same system was treated in the grand canonical ensemble: there
the density was $\rho=P/(1+P)$, which is the same equation of state solved for $P$. The pressure is a
real-analytic function of $\rho\in[0,1)$ and blows up at $\rho=1=\rho_{\mathrm{cp}}$, as
Proposition~\ref{SM:prop:kepler} predicts. There is no phase transition. That is not an accident of the
model.

\begin{theorem}[van Hove 1950]\label{SM:thm:vanhove}
Consider particles on a line interacting by a pair potential that has a hard core ($\phi(r)=+\infty$ for
$r<a$) and finite range ($\phi(r)=0$ for $r>R$), and is bounded on $[a,R]$. Then the thermodynamic
limit of the pressure is a real-analytic function of the temperature and the density: one-dimensional
systems of this kind have no phase transitions.
\end{theorem}

The proof uses a \emph{transfer operator}: by the hard core and the finite range, the partition function
is an iterated integral $\langle f,\mathcal T^ng\rangle$ of a positive integral operator, whose simple top
eigenvalue (Perron--Frobenius--Jentzsch) depends analytically on the parameters. Section~\ref{SM:sec:edwards}
has the same mechanism with a $2\times2$ matrix. (Cuesta and S\'anchez, J.~Stat.~Phys.\ 2004, discuss
which hypotheses are needed, and give one-dimensional models with transitions when they fail.)

\subsection*{Two and three dimensions: transitions seen only on the computer}

In $\R^3$ the hard-sphere gas is believed to \emph{freeze}. In 1957 Alder and Wainwright (molecular
dynamics) and Wood and Jacobson (Monte Carlo) published, back to back in J.~Chem.\ Phys.~27, the
first computer evidence of a fluid--solid transition. In the accepted picture (Hoover and Ree 1968) the
fluid is stable up to packing fraction about $0.494$, the close-packed crystal from about $0.545$, and
the two coexist in between. All of this is \emph{numerical}: that hard spheres crystallize, or have any
phase transition at all, is open in every dimension $d\ge2$. In $\R^2$ long-range positional order is
impossible (Mermin 1968 for smooth potentials; Richthammer 2016 for hard discs, see the gallery), and
what can survive is \emph{orientational} bond order, the hexatic phase. Bernard and Krauth (Phys.\ Rev.\
Lett.\ 107, 2011) found numerically that hard discs melt in two steps: a first-order liquid--hexatic
transition and a continuous hexatic--solid transition. Rigorous results are mostly at zero temperature,
where the problem becomes an optimization problem of the kind studied in this course (Theil's theorem in
the gallery). The first continuum model with finite-range interactions and a \emph{proved} phase
transition was the two-species Widom--Rowlinson model (Ruelle, Phys.\ Rev.\ Lett.\ 27, 1971), with hard
cores only between unlike particles.

%======================================================================
\section{Random packings}
\label{SM:sec:random}

\subsection{R\'enyi's parking problem}

Lecture~1 (Proposition~\ref{L1:prop:saturated-bound}) gave the easiest lower bound for packing
densities: a \emph{saturated} packing, one to which no ball can be added, has density at least
$2^{-d}$. The most natural random way of producing a saturated packing is to throw balls at random,
one at a time, and keep each one that does not overlap those already placed, until there is no room
left. This is \emph{random sequential adsorption} (RSA). In dimension one it is R\'enyi's \emph{parking
problem}: cars of length $1$ park at uniformly random free positions along a kerb of length $x$ until no
gap of length $\ge1$ remains.

Let $M(x)$ be the expected number of cars that park. Clearly $M(x)=0$ for $x<1$, and $M(x)=1$ for
$1\le x<2$. For $x\ge1$, condition on the position $t\in[0,x-1]$ of the first car, which is uniform. It
leaves two independent kerbs, of lengths $t$ and $x-1-t$:
\begin{equation}\label{SM:eq:renyi}
  M(x)=1+\frac1{x-1}\int_0^{x-1}\big(M(t)+M(x-1-t)\big)\,dt=1+\frac2{x-1}\int_0^{x-1}M(t)\,dt .
\end{equation}

\begin{theorem}[R\'enyi 1958]\label{SM:thm:renyi}
$M(x)/x\to m$ as $x\to\infty$, where
\[
  m=\int_0^\infty\exp\Big(-2\int_0^t\frac{1-e^{-u}}{u}\,du\Big)\,dt=0.7475979202\ldots
\]
\end{theorem}

\begin{proof}
Since the cars are disjoint subintervals of $[0,x]$, $0\le M(x)\le x$, so all Laplace transforms below
converge for $s>0$. Put $N(y)=M(y+1)$ for $y\ge0$ and multiply \eqref{SM:eq:renyi} by $y=x-1$:
\[
  yN(y)=y+2\int_0^yM(t)\,dt .
\]
Take Laplace transforms $\hat f(s)=\int_0^\infty e^{-sy}f(y)\,dy$. The transform of $yN(y)$ is
$-\hat N'(s)$, that of $\int_0^yM$ is $\hat M(s)/s$, and $\hat M(s)=e^{-s}\hat N(s)$ because $M$ vanishes on
$[0,1)$. So
\[
  \hat N'(s)+\frac{2e^{-s}}{s}\hat N(s)=-\frac1{s^2}.
\]
Solutions of the homogeneous equation are constant multiples of $\exp\big(2\int_s^\infty e^{-u}u^{-1}du\big)$,
which tends to a nonzero constant as $s\to\infty$, whereas $\hat N(s)\le\int_0^\infty e^{-sy}(y+1)\,dy\to0$. Hence
$\hat N$ is the unique decaying solution,
\[
  \hat N(s)=\int_s^\infty\frac{1}{t^2}\exp\Big(2\int_s^t\frac{e^{-u}}u\,du\Big)\,dt
  =\frac1{s^2}\int_s^\infty\exp\Big(-2\int_s^t\frac{1-e^{-u}}u\,du\Big)\,dt,
\]
where we wrote $e^{-u}=1-(1-e^{-u})$ and used $\exp(2\int_s^t du/u)=t^2/s^2$. As $s\to0^+$ the last
integral tends to $m$ by dominated convergence (for $t\ge1$ the integrand is at most
$\exp(-2\int_1^t\frac{1-e^{-u}}u du)=O(t^{-2})$). Hence $\hat M(s)=e^{-s}\hat N(s)\sim m/s^2$.

Now $U(x)=\int_0^xM$ is nondecreasing and $\int_0^\infty e^{-sx}dU(x)=\hat M(s)\sim m\,s^{-2}$, so Karamata's
Tauberian theorem (Feller, \emph{An Introduction to Probability Theory}, Vol.~II, \S XIII.5) gives
$U(x)\sim mx^2/2$. Finally \eqref{SM:eq:renyi} itself says $M(x)=1+2U(x-1)/(x-1)$, so $M(x)\sim mx$.
\end{proof}

\paragraph{Numerical check.} Quadrature gives $m=0.74759792025$ (R\'enyi, Publ.\ Math.\ Inst.\ Hung.\
Acad.\ Sci.\ 3, 1958). Twenty simulated kerbs of length $10^5$ gave
$0.74762\pm0.00012$. Solving \eqref{SM:eq:renyi} on a grid shows that $M(x)-mx$ settles down to
$m-1\approx-0.2524$ already for $x\approx10$. The density $m\approx0.748$ lies well above the saturation
bound $\frac12$ and well below~$1$.

\begin{figure}[ht]
\centering
\begin{tikzpicture}[scale=0.95]
  % 1D parking: first car splits the kerb
  \draw[thick] (0,0)--(7,0);
  \draw (0,-0.12)--(0,0.12) node[below=7pt,font=\small]{$0$};
  \draw (7,-0.12)--(7,0.12) node[below=7pt,font=\small]{$x$};
  \fill[flagcol!60] (2.6,0.05) rectangle (3.6,0.35);
  \draw (2.6,0.05) rectangle (3.6,0.35);
  \node[font=\scriptsize] at (3.1,0.6) {first car};
  \draw[<->] (0,-0.45)--node[below,font=\small]{$t$}(2.6,-0.45);
  \draw[<->] (3.6,-0.45)--node[below,font=\small]{$x-1-t$}(7,-0.45);
  \foreach \a in {0.3,1.45,4.1,5.2} {\fill[addcol!30] (\a,0.05) rectangle (\a+1,0.35); \draw (\a,0.05) rectangle (\a+1,0.35);}
  \node[font=\small,align=center] at (3.5,-1.6) {two independent smaller\\ parking problems};
  % 2D RSA picture
  \begin{scope}[xshift=8.0cm,yshift=-1.9cm,scale=0.5]
    \clip (0,0) rectangle (9,9);
    \foreach \p in {(5.63,8.07),(6.98,2.03),(2.70,7.86),(0.05,7.39),(9.05,7.39),(7.17,4.21),(2.73,2.51),(2.29,4.01),(4.54,4.98),(1.94,1.44),(5.51,0.40),(5.51,9.39),(0.32,4.63),(9.32,4.63),(4.20,8.26),(5.66,4.63),(4.47,2.23),(0.11,1.73),(9.11,1.73),(3.33,0.03),(3.33,9.03),(5.76,6.68),(3.25,5.38),(0.53,3.49),(6.09,1.36),(6.05,2.70),(7.87,5.96),(1.19,7.61),(5.13,3.39),(2.90,6.76),(7.86,3.10),(6.89,8.18),(3.77,7.34),(0.13,5.66),(9.13,5.66),(1.61,3.12),(6.68,5.23),(7.56,1.01),(0.27,0.18),(0.27,9.18),(9.27,0.18),(9.27,9.18),(1.69,5.10),(1.49,6.10),(7.76,7.17),(1.45,0.24),(1.45,9.24),(4.11,5.92),(3.40,3.78),(8.00,8.25),(-0.45,3.90),(8.55,3.90),(6.75,6.37),(6.25,3.79),(3.53,1.38),(4.72,6.71),(4.60,1.03),(5.32,5.65),(6.81,0.18),(6.81,9.19),(8.28,4.90),(8.03,2.01),(0.82,1.02),(0.95,2.32),(-0.14,2.70),(8.86,2.70),(2.48,0.57),(1.28,4.15),(8.21,0.25),(8.21,9.25),(0.60,6.55)} {\filldraw[fill=addcol!20,draw=addcol] \p circle (0.5);}
    \draw[thick] (0,0) rectangle (9,9);
  \end{scope}
\end{tikzpicture}
\caption{Left: the recursion behind R\'enyi's integral equation \eqref{SM:eq:renyi}. Right: a saturated
random sequential adsorption of $57$ discs of diameter $1$ in a periodic $9\times9$ square (computed for
this figure). Every point of the square is within distance $1$ of a centre, yet no two discs touch.}
\label{SM:fig:parking}
\end{figure}

In the plane, RSA of discs saturates at packing fraction about $0.547$ (Feder 1980). This is a
numerical value; Zhang and Torquato (Phys.\ Rev.\ E 88, 2013) give $0.5470735(28)$, and in $\R^3$
$0.3841307(21)$. (A quick simulation for this interlude, ten saturated periodic configurations of about
$1740$ discs each, gave $0.5468\pm0.0009$.) No closed form is known in any dimension $d\ge2$.

Notice what RSA does \emph{not} produce. In a saturated RSA configuration, with probability one no two
discs touch: the contact graph of Lectures~10--11 is empty. RSA packings are saturated but not
jammed; a small shake would let every disc move. To get mechanically stable random packings one has to
compress.

\subsection{Random close packing and the maximally random jammed state}

Pour ball bearings into a container and shake it: the packing fraction settles near $0.64$. This
\emph{random close packing} (RCP) value goes back to the experiments of Bernal and of Scott (both in
\emph{Nature}, 1960); Scott and Kilgour (1969) measured about $0.637$. It is well below the crystal value
$0.7405$, and remarkably reproducible.

Torquato, Truskett and Debenedetti (\emph{Is random close packing of spheres well defined?}, Phys.\
Rev.\ Lett.\ 84, 2000) argued that it is not: by compressing faster or slower one trades a little
disorder for a little density, with no canonical place to stop. They proposed instead the
\emph{maximally random jammed} (MRJ) state, the jammed packing minimizing a chosen order parameter,
again at about $0.64$ (numerical). This needs a definition of ``jammed''; Torquato and Stillinger
(J.~Phys.\ Chem.~B 105, 2001) give a hierarchy. For a packing in a container with rigid walls:
\begin{itemize}[nosep]
\item \emph{locally jammed}: no single ball can be moved while all the others are held fixed;
\item \emph{collectively jammed}: locally jammed, and no subset of balls can be moved
  collectively, the walls being fixed;
\item \emph{strictly jammed}: collectively jammed, also under small volume-nonincreasing
  deformations of the walls (for periodic packings: of the lattice).
\end{itemize}
For instance, with suitable walls, the simple cubic packing is collectively but not strictly jammed (shearing the box lets it
slide), FCC is strictly jammed, and the honeycomb disc packing is locally but not collectively jammed.
(In a torus there are no walls to stop a row, and the square and simple cubic packings are not even
collectively jammed; see Figure~\ref{SM:fig:maxwell}.)
To first order, a ball with the others held fixed is stuck exactly when the directions to its contacts
do not all lie in a closed hemisphere. In particular a locally jammed ball has at least $d+1$ contacts.

\subsection{Jamming, force balance and Maxwell counting}

Collective jamming is the version of Lectures~10--11 that sees all balls at once. Take $N$ balls with
centres $x_i$ and radii $r_i$ in a flat torus $\R^d/\Lambda$ (periodic boundary conditions), with contact
set $E=\{\{i,j\}:\abs{x_i-x_j}=r_i+r_j\}$, and let $u_{ij}=(x_j-x_i)/\abs{x_j-x_i}$. A velocity
$v=(v_1,\dots,v_N)\in(\R^d)^N$ changes the contact distances at the rates
\[
  (Av)_{ij}=\ip{u_{ij}}{v_j-v_i},\qquad\{i,j\}\in E.
\]
The linear map $A\colon(\R^d)^N\to\R^E$ is the \emph{rigidity matrix} of the contact graph. To first order a
motion keeps the balls disjoint if and only if $Av\ge0$ componentwise. Translations $v_i=c$ always
satisfy $Av=0$; call them \emph{trivial}.

\begin{definition}\label{SM:def:jam}
The packing is \emph{first-order jammed} if every $v$ with $Av\ge0$ is trivial.
\end{definition}

For sphere packings, first-order jamming is equivalent to collective jamming (with the lattice fixed):
a sphere packing is rigid if and only if it is infinitesimally rigid. This is a theorem of Connelly (1988)
for packings in a container, extended to periodic packings by Donev, Torquato, Stillinger and Connelly
(J.~Comput.\ Phys.\ 197, 2004). It is special to spheres; it fails for ellipsoids.

\begin{proposition}\label{SM:prop:jam}
A packing is first-order jammed if and only if
\begin{enumerate}[nosep]
\item $\ker A$ consists of the trivial motions (the contact graph is infinitesimally rigid as a bar
  framework), and
\item there are \emph{strictly positive} weights $w_{ij}>0$ on all contacts with
  $\sum_jw_{ij}u_{ij}=0$ at every ball $i$ (a positive self-stress: the forces balance at every ball).
\end{enumerate}
\end{proposition}

\begin{proof}
If (1) and (2) hold and $Av\ge0$, then
\[
  0=\sum_i\Big\langle v_i,\sum_jw_{ij}u_{ij}\Big\rangle=-\sum_{\{i,j\}\in E}w_{ij}\ip{u_{ij}}{v_j-v_i}=-\ip{w}{Av},
\]
a sum of non-positive terms, so $Av=0$, and $v$ is trivial by (1). Conversely, (1) is necessary because
$Av=0$ implies $Av\ge0$. For (2), Stiemke's alternative (a relative of Gordan's alternative used in the
proof of Theorem~\ref{L11:thm:forcebalance}) says that either there is $v$ with $Av\ge0$ and $Av\ne0$, or
there is $w>0$ with $A^Tw=0$. In a first-order jammed packing the first option is excluded, and $A^Tw=0$
is exactly the balance condition.
\end{proof}

So the jammed packings are the force-balanced contact graphs of Lecture~10 with two upgrades: the
weights are positive on \emph{every} contact, and the graph is rigid. (Compare
Corollary~\ref{L10:cor:topchange}, where critical configurations of the min-distance function only
needed weights $\ge0$, not all zero.) The rigidity condition already forces a lot of contacts.

\begin{corollary}[Maxwell counting]\label{SM:cor:maxwell}
A first-order jammed packing of $N$ balls in a $d$-dimensional torus has at least $dN-d+1$ contacts, so the
mean number of contacts per ball satisfies
\[
  z=\frac{2\abs E}{N}\ \ge\ 2d-\frac{2(d-1)}{N}.
\]
\end{corollary}

\begin{proof}
By (1), $\ker A$ has dimension $d$, so $\operatorname{rank}A=dN-d$. By (2), $A^T$ has a nonzero kernel, so
$\abs E\ge\operatorname{rank}A+1=dN-d+1$.
\end{proof}

This is Maxwell's 1864 count for frames. Conversely, for generic radii each contact is one equation
$\abs{x_i-x_j}=r_i+r_j$ on $dN-d$ free coordinates plus one free parameter, the common scale of the radii
at which the packing jams, and generic equations should not be overdetermined.
Together the counts give $z=2d$: jammed frictionless packings should be \emph{isostatic}. The upper bound
is a theorem in the plane (Connelly--Gortler--Solomonides--Yampolskaya, gallery) and open for $d\ge3$, but
it holds very precisely in simulations. O'Hern, Silbert, Liu and Nagel (Phys.\ Rev.~E 68, 2003), compressing
soft repulsive spheres to the threshold ``point~J'' where they first jam, found $z=3.98$ in $2$d and $5.98$ in
$3$d (both $\pm0.02$, rattlers removed) and a $3$d monodisperse threshold $\phi^*=0.639\pm0.001$,
essentially the RCP value.

\begin{figure}[ht]
\centering
\begin{tikzpicture}[scale=0.55]
  \foreach \i in {0,...,5} \foreach \j in {0,1,3} {\draw (\i,\j) circle (0.5);}
  \foreach \i in {0,...,5} {\filldraw[fill=addcol!20,draw=addcol] (\i+0.35,2) circle (0.5);}
  \draw[-{Stealth},thick,flagcol] (5.6,2)--(6.6,2);
  \draw[dashed,gray] (-0.6,-0.6) rectangle (5.6,3.6);
  \node[right,font=\small,align=left] at (7.0,1.5) {square lattice in a torus: $z=4=2d$,\\
    but the shaded row slides; to first order\\ $Av=0$ for this motion, so it is not jammed};
\end{tikzpicture}
\caption{Maxwell's count is necessary, not sufficient. In the square packing every disc has four
contacts, yet a whole row can slide along itself: its vertical contacts open only to second order.}
\label{SM:fig:maxwell}
\end{figure}


%======================================================================
\section{Geometric frustration}
\label{SM:sec:frustration}

Why does a dense liquid not simply crystallize? In 1952 F.~C.~Frank (\emph{Supercooling of liquids},
Proc.\ R.\ Soc.\ Lond.\ A 215) pointed out that the question has a \emph{local} answer. Put twelve atoms
around a thirteenth. There are three natural candidates: the cuboctahedron (the FCC neighbourhood), its
twisted relative from HCP, and the icosahedron. Lecture~11 met the first and the last, which the jitterbug
connects. For
Lennard-Jones atoms Frank computed that the icosahedral cluster has the lowest energy of the three: it
has $30$ near-neighbour pairs among the outer twelve, against $24$. But icosahedral symmetry is
incompatible with a lattice, so a liquid of locally icosahedral clusters must rearrange substantially
before it can crystallize, and it can be supercooled. This is the first instance of \emph{geometric
frustration}: the locally preferred structure cannot be extended to fill space.

For hard spheres the same tension appears as a statement about densities. The densest \emph{local}
arrangement around one ball beats the densest global one. Fejes T\'oth conjectured, and Hales and
McLaughlin proved (J.~Amer.\ Math.\ Soc.\ 23, 2010), that in any packing of unit balls in $\R^3$ every
Voronoi cell (Lecture~1, Definition~\ref{L1:def:voronoi}) has volume at least that of the regular
dodecahedron of inradius $1$. That cell would give density
\[
  \frac{\vol\Ball^3}{\vol(\text{dodecahedron of inradius }1)}\approx0.7547\ >\ \frac{\pi}{\sqrt{18}}\approx0.7405 .
\]
The dodecahedral cell is the Voronoi cell of a ball surrounded by an icosahedron of twelve touching balls,
and regular dodecahedra do not tile space. The local optimum is not the global one.

Frank's icosahedron is often described as twenty regular tetrahedra, the densest arrangement of four
balls, glued around the centre. That description is slightly wrong, and the error is the heart of the
matter.

% Section for Interlude E: why 20 regular tetrahedra do not make a regular icosahedron.
% Written to follow W. Kusner's argument. Labels prefixed TET:.

\section{Twenty tetrahedra do not make an icosahedron}
\label{TET:sec}

Frank's icosahedral cluster (and the jitterbug of Lecture~11) rests on a small geometric fact: a regular
icosahedron is \emph{not} made of $20$ regular tetrahedra. Cone the icosahedron from its centre and
you get $20$ congruent tetrahedra, but they are slightly flattened. The edges from the centre are a
little \emph{shorter} than the edges of the icosahedron:
\[
  \frac{\text{circumradius}}{\text{edge}}=\sin\frac{2\pi}{5}\approx0.951<1 .
\]
In terms of balls: twelve unit balls placed at the vertices of an icosahedron around a central unit
ball can all touch the central ball, but then they do not touch each other. This slack is why the
twelve outer balls can roll around (the kissing configuration in $\R^3$ is not rigid), and it is the
root of the Newton--Gregory dispute over whether a thirteenth ball fits.

The usual proof is a computation, for instance with the dihedral angle $\arccos\frac13$ of the regular
tetrahedron. The following argument needs no computation. It uses only the plane fact that \emph{six}
equal circles fit exactly around a seventh, together with the triangle inequality on the sphere.

\begin{proposition}\label{TET:prop}
A regular icosahedron is not the union of $20$ regular tetrahedra sharing a vertex at its centre.
\end{proposition}

\paragraph{The argument.} Suppose it were, with edge length $a$. Put balls of radius $a/2$ at the centre
and at the $12$ vertices. Then the central ball touches all twelve, and each vertex ball touches its
five neighbours. Project the twelve contact points radially onto a sphere about the centre, and
measure distances by angle. Neighbouring vertices are then exactly $60^\circ$ apart, because the centre
and two neighbours form an equilateral triangle.

Stand the icosahedron on a vertex: a north pole $N$, the antipodal vertex $S$, and two
\emph{tropical rings} of five vertices each. The upper ring consists of the neighbours of $N$, the
lower ring of the neighbours of $S$, and the rings are staggered by $36^\circ$ of longitude. The equatorial
band between the rings is made of ten triangles (Figure~\ref{TET:fig:band}, left). Every vertex of the
upper ring is $60^\circ$ from $N$, so the upper ring lies at latitude $30^\circ$\,N. Likewise the lower ring
lies at latitude $30^\circ$\,S.

Now \emph{rotate the lower ring} by $36^\circ$ about the axis $NS$, keeping its latitude. Each lower vertex
$\ell$ now sits directly below an upper vertex $u$, and the ten equatorial triangles have become five
equatorial quadrilaterals, the ``squares'' (Figure~\ref{TET:fig:band}, right). Each of the five
meridians through $N$ and $S$ now carries four of our points,
\[
  N,\ u,\ \ell,\ S\qquad\text{at angles}\qquad 0^\circ,\ 60^\circ,\ 120^\circ,\ 180^\circ
\]
along the meridian, since $u$ is $60^\circ$ from $N$, $\ell$ is $60^\circ$ from $S$, and $N$ and $S$ are
$180^\circ$ apart. The meridian is a great circle, and $60^\circ$ steps are exactly the spacing of a
\emph{ring of six} balls around the central one, as in the hexagonal packing
(Figure~\ref{TET:fig:meridian}). So after the rotation, $u$ and $\ell$ touch: the five meridians have
become five half-rings of six, one great circle each.

But before the rotation $u$ and $\ell$ also touched, since they were neighbours in the icosahedron, and
in between we only slid $\ell$ sideways along its latitude circle. The closest point of a latitude
circle to $u$ is the one on $u$'s meridian, so sliding $\ell$ off the meridian moves it strictly
\emph{farther} from $u$. In other words, the deformation from ten equatorial triangles to five
equatorial squares, made without changing any latitude, shortened the distance from $u$ to $\ell$. Yet
it is $60^\circ$ at both ends. That is a contradiction. \qed

\medskip
In one line, the same argument is the triangle inequality on the sphere:
\[
  d(u,\ell)\ \ge\ d(N,\ell)-d(N,u)\ =\ (180^\circ-60^\circ)-60^\circ\ =\ 60^\circ,
\]
with equality only if $u$ lies on the arc from $N$ to $\ell$, that is, on $\ell$'s meridian. Each upper
vertex $u$ has two lower neighbours, and they cannot both lie on $u$'s meridian, so for at least one of
them $d(u,\ell)>60^\circ$. (In the icosahedron both are off the meridian, by $\pm36^\circ$.) The ring of six
is the picture of the equality case.

\begin{figure}[ht]
\centering
\begin{tikzpicture}[scale=0.95]
  % left: antiprism band (10 triangles), right: prism band (5 squares)
  \begin{scope}
    \foreach \k in {0,...,5} {\coordinate (U\k) at (\k,1); }
    \foreach \k in {0,...,4} {\coordinate (L\k) at (\k+0.5,0); }
    \coordinate (Lm) at (-0.5,0);
    \draw[thick] (U0)--(U5); \draw[thick] (Lm)--(L4);
    \foreach \k [evaluate=\k as \j using int(\k+1)] in {0,...,4} {\draw (U\k)--(L\k)--(U\j);}
    \draw (U0)--(Lm);
    \foreach \k in {0,...,5} \fill (U\k) circle (0.06);
    \foreach \k in {0,...,4} \fill[flagcol] (L\k) circle (0.06);
    \node[font=\small] at (2.5,1.55) {upper ring, latitude $30^\circ$\,N};
    \node[font=\small] at (2.5,-0.55) {lower ring, latitude $30^\circ$\,S};
    \node[font=\small,gray] at (2.5,-1.1) {ten triangles (icosahedron)};
  \end{scope}
  \draw[-{Stealth},thick] (5.8,0.5)--(6.9,0.5) node[midway,above,font=\scriptsize]{rotate};
  \node[font=\scriptsize] at (6.35,0.2) {lower ring};
  \begin{scope}[xshift=7.4cm]
    \foreach \k in {0,...,5} {\coordinate (U\k) at (\k,1); \coordinate (L\k) at (\k,0);}
    \draw[thick] (U0)--(U5); \draw[thick] (L0)--(L5);
    \foreach \k in {0,...,5} \draw (U\k)--(L\k);
    \foreach \k in {0,...,5} {\fill (U\k) circle (0.06); \fill[flagcol] (L\k) circle (0.06);}
    \node[font=\small,gray] at (2.5,-1.1) {five squares, same latitudes};
  \end{scope}
\end{tikzpicture}
\caption{The equatorial band, unrolled (the left and right ends are identified; longitudes in units of
$72^\circ$). Rotating the lower ring by $36^\circ$ turns the ten equatorial triangles into five
quadrilaterals without changing any latitude.}
\label{TET:fig:band}
\end{figure}

\begin{figure}[ht]
\centering
\begin{tikzpicture}[scale=0.62]
  \fill[gray!15] (0,0) circle (1); \draw (0,0) circle (1); \node[font=\scriptsize] at (0,0) {$O$};
  \foreach \a/\lab in {90/N,30/u,-30/\ell,-90/S}
    {\filldraw[fill=addcol!15,draw=addcol,thick] (\a:2) circle (1); \node[font=\small] at (\a:2) {$\lab$};}
  \foreach \a in {150,210} {\draw[dashed,gray] (\a:2) circle (1);}
  \draw[gray,thin] (0,0) circle (2);
  \foreach \a in {90,30,-30,-90} \draw[gray,thin] (0,0)--(\a:2);
  \node[font=\scriptsize] at (60:0.65) {$60^\circ$};
  \node[font=\scriptsize] at (0:0.7) {$60^\circ$};
  \node[font=\scriptsize] at (-60:0.65) {$60^\circ$};
\end{tikzpicture}
\caption{After the rotation, each meridian is half of a ring of six balls around the central ball: $N$,
$u$, $\ell$, $S$ are consecutive at $60^\circ$ steps. The dashed balls complete the ring.}
\label{TET:fig:meridian}
\end{figure}

\begin{remark}[What the argument shows, quantitatively]
In the hypothetical configuration with rings at latitudes $\pm30^\circ$, an upper and a lower vertex
whose longitudes differ by $36^\circ$ are at angular distance $\arccos\big(\tfrac34\cos36^\circ-\tfrac14\big)
\approx69.1^\circ$, not $60^\circ$. In the true icosahedron the rings sit lower, at latitudes
$\pm\arctan\frac12\approx\pm26.6^\circ$, every edge subtends $\arctan2\approx63.4^\circ$ at the centre, and
the gaps between neighbouring outer balls are what let the twelve balls roll.
\end{remark}

\begin{remark}[Two related classical facts]
(1) Regular tetrahedra do not even fill the space around an \emph{edge}: the dihedral angle is
$\arccos\frac13\approx70.53^\circ$, and five such angles leave a gap of about $7.36^\circ$. (Aristotle
claimed that regular tetrahedra tile space; see Lagarias--Zong, \emph{Mysteries in packing regular
tetrahedra}, Notices AMS 59 (2012).) The same angle appears here: five points at $60^\circ$ from $N$ and
$60^\circ$ from each other would need angles of $\arccos\frac13$ at $N$, and five of these do not close up.
(2) In curved space the frustration disappears. In the $3$-sphere, $600$ regular tetrahedra (with
the right curvature) fit together exactly, five around each edge, forming the regular polytope
$\{3,3,5\}$ with $120$ vertices, each surrounded by a regular icosahedron. This is the ``ideal''
icosahedral order of Frank's liquids; the flat-space packing problem has to compromise.
\end{remark}

\begin{remark}[Source]
\unclear{Wording of the attribution is a placeholder for the author to edit.}
This argument came out of a conversation between the author and T.~C.~Hales.
\end{remark}


%======================================================================
\section{Ideal order in curved space}
\label{SM:sec:curved}

The previous section measured the frustration: five regular tetrahedra around an edge leave a gap of
about $7.36^\circ$. As noted there, the gap closes in positive curvature. The dihedral angle of a regular
spherical tetrahedron grows with its size from the flat value $\arccos\frac13\approx70.53^\circ$, and for
edge $36^\circ$ on the unit $\Sph^3$ it is exactly $72^\circ=360^\circ/5$. Six hundred such tetrahedra are the cells of the
regular polytope $\{3,3,5\}$ (Coxeter, \emph{Regular Polytopes}), whose $120$ vertices are at angular
distance at least $36^\circ$, each with $12$ nearest neighbours forming a regular icosahedron. Balls of
angular radius $18^\circ$ at these points realize Frank's local structure \emph{everywhere}. By a theorem of
Cohn and Kumar (gallery) these $120$ points even minimize every completely monotonic pair potential on
$\Sph^3$: the ideal order is a zero-temperature ground state, though in curved space.

This is the starting point of the curved-space approach to liquids and glasses of Sadoc, Mosseri and
others, and of Nelson (\emph{Order, frustration, and defects in liquids and glasses}, Phys.\ Rev.~B 28,
1983): a dense flat liquid is the ideal $\{3,3,5\}$ structure ``flattened'', with the curvature
concentrated on line defects (disclinations) along which six rather than five tetrahedra share an edge.
See Sadoc and Mosseri, \emph{Geometrical Frustration} (Cambridge, 1999).

%======================================================================
\section{Edwards entropy}
\label{SM:sec:edwards}

\subsection{The ansatz and its tests}

A pile of sand is not at thermal equilibrium; $k_BT$ is negligible next to the energy needed to lift a
grain. Yet tapped or sheared granular packings settle into reproducible states, and Edwards and
Oakeshott (\emph{Theory of powders}, Physica~A 157, 1989) proposed a statistical mechanics for them in
which the \emph{volume} plays the role of energy. This is Definition~\ref{L11:def:edwards}: all jammed
(mechanically stable) states of the same volume are equally probable. Writing $\Omega(V)$ for the number
(or a suitable measure) of jammed states of volume $V$, the dictionary is
\begin{center}
\renewcommand{\arraystretch}{1.2}
\begin{tabular}{ll}
\toprule
thermal statistical mechanics & Edwards' granular ensemble\\
\midrule
energy $E$ & volume $V$\\
all states of energy $E$ & jammed states of volume $V$\\
entropy $S(E)=\ln\Omega(E)$ & Edwards entropy $S(V)=\ln\Omega(V)$\\
temperature $T=\partial E/\partial S$ & compactivity $X=\partial V/\partial S$\\
Boltzmann weight $e^{-E/T}$ & weight $e^{-V/X}$\\
\bottomrule
\end{tabular}
\end{center}
The compactivity measures how much more compact the packing could be. As $X\to0^+$ the weight
concentrates on the densest jammed states, and as $X\to\infty$ on the most numerous ones. Compare
Section~\ref{SM:sec:hard}: for the thermal hard-sphere gas $\partial S/\partial V=\beta P$, while for
Edwards $\partial S/\partial V=1/X$.

In the continuum, jammed packings are isolated points of configuration space, so ``equal
probability'' needs a measure, and a real protocol samples a jammed state with probability equal to the
volume of its \emph{basin of attraction}, as in the $\rho$-evolution of Lecture~11. Edwards' ansatz then
says that basins of equal-volume states are equally large; as noted in Lecture~11, small disc packings
violate this. The most careful test so far (Martiniani, Schrenk, Ramola, Chakraborty and Frenkel,
\emph{Nature Physics} 13, 2017) computed basin volumes by thermodynamic integration for $N=64$
polydisperse discs in $2$d (hard cores with soft repulsive shells), with $N=32$ to $128$ for finite-size
scaling. Above jamming the basin volumes are broadly distributed and the ansatz fails; the distribution
narrows towards the unjamming point, extrapolated to $\phi^*=0.841(3)$, where all packings become
equally probable (numerically). Baule, Morone, Herrmann and Makse (Rev.\ Mod.\ Phys.\ 90, 2018) review
the theory and its tests.

Where Edwards' ensemble can be written down \emph{exactly}, its predictions can be checked against the
ansatz-free dynamics. The channel of Lecture~11 is the simplest such model, and we now solve it.

\subsection{The channel model solved}
\label{SM:sec:channel}

Recall the setting of \S\ref{L11:sec:channel}. Discs of radius $1$ lie in a channel of width $w=2+h$
with $0<h<\sqrt3$, so that each disc touches at most its two neighbours. Two consecutive discs in
contact at opposite walls are horizontally $L=\sqrt{4-h^2}\in(1,2)$ apart, and two at the same wall are
$2$ apart. A jammed configuration of $n=\ell+1$ discs is recorded by its contact word
$\sigma\in\{0,1\}^\ell$ with no two adjacent $1$'s ($1$ for same wall, $0$ for opposite walls). A word
with $k$ ones has length
\[
  t(\sigma)=2k+(\ell-k)L+2=\ell L+(2-L)k+2,
\]
and the number of such words is
\[
  \Omega(\ell,k)=\binom{\ell-k+1}{k},
\]
or $\binom{\ell-k-1}{k}$ if the end discs are restricted, as flagged in Lecture~11.
We take the volume to be the channel length, $V=t$ (multiply by $w$ for the area). Since $t$ is an affine
function of $k$, fixing the volume is the same as fixing the number $k$ of defects. Write
$x=k/\ell\in[0,\frac12]$ for the defect density, and use $\eta$ for the packing fraction (the letter
$\varphi$ is reserved for the golden ratio).

\begin{proposition}[Edwards entropy of the channel]\label{SM:prop:entropy}
As $\ell\to\infty$ with $k/\ell\to x\in[0,\frac12]$,
\[
  \frac1\ell\ln\Omega(\ell,k)\ \longrightarrow\ s(x)=(1-x)\ln(1-x)-x\ln x-(1-2x)\ln(1-2x),
\]
for both counts. The function $s$ is strictly concave on $[0,\frac12]$, vanishes at both ends, and has
its maximum $s(x^*)=\ln\varphi$ at $x^*=\frac{5-\sqrt5}{10}$. The volume per contact is
$v(x)=L+(2-L)x$ and the packing fraction is
\[
  \eta(x)=\frac{\pi}{w\,v(x)}=\frac{\pi}{w\,(L+(2-L)x)} .
\]
\end{proposition}

\begin{proof}
Stirling's formula in the form $\ln\binom ab=a\,H\!\big(\tfrac ba\big)+O(\ln a)$, with
$H(p)=-p\ln p-(1-p)\ln(1-p)$, applied with $a=\ell(1-x)+O(1)$ and $b=\ell x$, gives the limit
$(1-x)H\big(\frac x{1-x}\big)$. Expanding, with $1-\frac x{1-x}=\frac{1-2x}{1-x}$, this is $s(x)$. The
shift $O(1)$ in $a$ does not change the limit, so the two counts agree to leading order. Next,
\[
  s'(x)=\ln\frac{(1-2x)^2}{x(1-x)},\qquad s''(x)=\frac1{1-x}-\frac1x-\frac4{1-2x}<0,
\]
since $\frac1{1-x}\le2<\frac4{1-2x}$ on $[0,\frac12)$. The critical point solves $(1-2x)^2=x(1-x)$, that is
$5x^2-5x+1=0$, as in Lecture~11. At $x^*$ one has $1-x^*=\frac{5+\sqrt5}{10}=\frac{\varphi}{\sqrt5}$,
$x^*=\frac{1}{\sqrt5\,\varphi}$ and $1-2x^*=\frac{1}{\sqrt5}$, so that
$s(x^*)=\frac\varphi{\sqrt5}\ln\frac\varphi{\sqrt5}-\frac1{\sqrt5\varphi}\ln\frac1{\sqrt5\varphi}-\frac1{\sqrt5}\ln\frac1{\sqrt5}$.
Since $\frac\varphi{\sqrt5}-\frac1{\sqrt5\varphi}-\frac1{\sqrt5}=0$, the terms in $\ln\sqrt5$ cancel and what is
left is $\big(\frac\varphi{\sqrt5}+\frac1{\sqrt5\varphi}\big)\ln\varphi=\ln\varphi$, using
$\varphi+\varphi^{-1}=\sqrt5$. The formula for $v$ is $t/\ell$ up to $O(1/\ell)$, and $\eta$ is the area
$\pi n$ of the discs divided by the area $wt$ of the channel, in the limit.
\end{proof}

The densest jammed state is the zigzag, $x=0$, with $\eta_{\max}=\pi/(wL)$; the loosest is
$x=\frac12$, with $\eta_{\min}=2\pi/(w(L+2))$. The zigzag word $00\cdots0$ is unique, the loosest word
$1010\cdots1$ is unique for odd $\ell$ (for even $\ell$ there are $\frac\ell2+1$ words with $\frac\ell2$
ones), and both have Edwards entropy zero. For $h=1.2$, the channel of Figure~\ref{L11:fig:channel}, one has $L=1.6$, $w=3.2$ and
\[
  \eta_{\max}\approx0.6136,\qquad \eta(x^*)\approx0.5739,\qquad \eta_{\min}\approx0.5454 .
\]

\paragraph{Compactivity and the equation of state.} Changing $k$ by one changes $V$ by $2-L$ and $\ln\Omega$
by $s'(x)+o(1)$, so
\begin{equation}\label{SM:eq:compactivity}
  \frac1X=\frac{\partial S}{\partial V}=\frac{s'(x)}{2-L}=\frac1{2-L}\ln\frac{(1-2x)^2}{x(1-x)} .
\end{equation}

\begin{proposition}[Equation of state]\label{SM:prop:eos}
For $X\in(-\infty,0)\cup(0,\infty]$, the defect density at compactivity $X$ is
\[
  x(X)=\frac12\Big(1-\frac1{\sqrt{1+4e^{-(2-L)/X}}}\Big),
\]
and $\eta(X)=\pi/\big(w(L+(2-L)x(X))\big)$. In particular:
\begin{itemize}[nosep]
\item as $X\to0^+$, $x\sim e^{-(2-L)/X}\to0$ and $\eta\to\eta_{\max}$ (densest);
\item $X=\pm\infty$ corresponds to $x=x^*$, the most numerous states;
\item the states looser than typical, $x^*<x<\frac12$, have \emph{negative} compactivity, and as $X\to0^-$,
  $x\to\frac12$ and $\eta\to\eta_{\min}$ (loosest).
\end{itemize}
\end{proposition}

\begin{proof}
Put $y=e^{(2-L)/X}$. By \eqref{SM:eq:compactivity}, $(1-2x)^2=y\,x(1-x)$, that is
$(4+y)x^2-(4+y)x+1=0$, whose root in $[0,\frac12]$ is $x=\frac12\big(1-\sqrt{y/(y+4)}\big)$, which is the
formula. The limits follow by inspection: $y\to\infty$, $y=1$, and $y\to0$ respectively. The sign of
$X$ is the sign of $s'(x)$, which changes at $x^*$.
\end{proof}

The small-$X$ behaviour $x\approx e^{-(2-L)/X}$ is an Arrhenius law: a defect costs ``volume'' $2-L$,
just as a thermal excitation costs energy. The negative branch is the analogue of negative temperatures
for spin systems, whose energy is bounded above. Here the jammed volume is bounded above, and past the
entropy maximum adding volume \emph{decreases} the number of states. For $h=1.2$, for example,
$X=0.25,1,-1,-0.25$ give $x\approx0.128,0.239,0.311,0.390$ and $\eta\approx0.595,0.579,0.569,0.559$.

\begin{figure}[ht]
\centering
\begin{tikzpicture}[scale=1.0]
  \draw[->] (0,0)--(7.6,0) node[below right,font=\small]{$\eta$};
  \draw[->] (0,0)--(0,4.5) node[left,font=\small]{$s$};
  \foreach \e/\x in {0.55/0.95,0.56/1.9,0.57/2.85,0.58/3.8,0.59/4.75,0.60/5.7,0.61/6.65}
    {\draw (\x,0)--(\x,-0.08) node[below,font=\scriptsize]{$\e$};}
  \foreach \v in {0.1,0.2,0.3,0.4} {\draw (0,8*\v)--(-0.08,8*\v) node[left,font=\scriptsize]{$\v$};}
  \draw[dashed,gray] (0,3.8497)--(3.224,3.8497) node[pos=0,left,font=\scriptsize,black]{$\ln\varphi$};
  \draw[dashed,gray] (3.224,0)--(3.224,3.8497);
  \draw[very thick,addcol] plot[smooth] coordinates {(6.991,0.000) (6.955,0.140) (6.918,0.252) (6.882,0.353) (6.846,0.447) (6.701,0.781) (6.557,1.071) (6.414,1.330) (6.272,1.567) (6.130,1.785) (5.989,1.987) (5.848,2.174) (5.709,2.349) (5.570,2.512) (5.431,2.663) (5.293,2.804) (5.156,2.935) (5.020,3.057) (4.884,3.169) (4.749,3.272) (4.615,3.367) (4.481,3.452) (4.348,3.530) (4.215,3.599) (4.084,3.660) (3.952,3.712) (3.822,3.756) (3.692,3.792) (3.562,3.819) (3.434,3.838) (3.305,3.848) (3.178,3.849) (3.051,3.841) (2.924,3.824) (2.799,3.798) (2.673,3.761) (2.549,3.715) (2.425,3.657) (2.301,3.589) (2.178,3.509) (2.056,3.416) (1.934,3.310) (1.813,3.191) (1.692,3.055) (1.572,2.903) (1.452,2.733) (1.333,2.542) (1.215,2.328) (1.097,2.086) (0.979,1.812) (0.862,1.497) (0.746,1.128) (0.688,0.915) (0.659,0.799) (0.630,0.675) (0.601,0.541) (0.572,0.393) (0.543,0.224) (0.514,0.000)};
  \fill (6.991,0) circle (0.06) node[above right,font=\scriptsize,align=left]{zigzag\\ $X\to0^+$};
  \fill (0.514,0) circle (0.06) node[right,xshift=3pt,yshift=6pt,font=\scriptsize,align=left]{$0101\cdots$, $X\to0^-$};
  \fill (3.224,3.8497) circle (0.06) node[above,font=\scriptsize]{$x^*=\frac{5-\sqrt5}{10}$, $X=\pm\infty$};
  \node[font=\small,flagcol] at (5.6,3.6) {$X>0$};
  \node[font=\small,flagcol] at (1.0,3.6) {$X<0$};
\end{tikzpicture}
\caption{Edwards entropy per contact of the channel model as a function of the packing fraction, for
$h=1.2$ ($L=1.6$, $w=3.2$); computed from Proposition~\ref{SM:prop:entropy}. Since $V$ decreases as $\eta$
increases, the compactivity $X=\partial V/\partial S$ is positive on the dense side of the maximum and
negative on the loose side.}
\label{SM:fig:entropy}
\end{figure}

\paragraph{The canonical ensemble is a transfer matrix.} Instead of fixing $V$, weight every jammed
state by $e^{-V/X}$:
\[
  Z_\ell(X)=\sum_\sigma e^{-t(\sigma)/X}=e^{-2/X}\sum_\sigma a^{\#0(\sigma)}\,b^{\#1(\sigma)},\qquad a=e^{-L/X},\ b=e^{-2/X},
\]
the sum running over words of length $\ell$ with no $11$. Keep track of the last letter. If
$(c_\ell(0),c_\ell(1))$ are the weighted sums over words ending in $0$ and in $1$, then
$(c_1(0),c_1(1))=(a,b)$ and
\[
  \big(c_{\ell+1}(0),c_{\ell+1}(1)\big)=\big(c_\ell(0),c_\ell(1)\big)\,T,\qquad
  T=\begin{pmatrix}a&b\\a&0\end{pmatrix},
\]
since after a $0$ we may write either letter and after a $1$ only a $0$. So we have
$Z_\ell(X)=e^{-2/X}(a,b)\,T^{\ell-1}(1,1)^T$. Since $T\ge0$ and $T^2>0$, Perron--Frobenius gives
\[
  \frac1\ell\ln Z_\ell(X)\longrightarrow\ln\lambda(X),\qquad \lambda(X)=\frac{a+\sqrt{a^2+4ab}}2 ,
\]
the larger root of $\lambda^2-a\lambda-ab=0$. This is analytic in $X\ne0$: as in van Hove's theorem, the
one-dimensional channel has no Edwards phase transition.

\begin{proposition}[Equivalence of ensembles]\label{SM:prop:canon}
The mean defect density in the canonical ensemble at compactivity $X$,
$\lim_\ell\langle k\rangle/\ell=b\,\partial_b\ln\lambda$ (with $a$ held fixed), equals the microcanonical
$x(X)$ of Proposition~\ref{SM:prop:eos}. Moreover
\[
  \ln\lambda(X)=\max_{0\le x\le1/2}\Big(s(x)-\frac{v(x)}X\Big),
\]
so $-X\ln\lambda$ is the Legendre transform of the entropy (an ``Edwards free energy'' per contact).
\end{proposition}

\begin{proof}
Write $\lambda=a\,g(u)$ with $u=b/a=e^{-(2-L)/X}$ and $g(u)=\frac12(1+r)$, $r=\sqrt{1+4u}$. Then
\[
  b\,\partial_b\ln\lambda=u\,\frac{g'(u)}{g(u)}=\frac{u\cdot 2/r}{1+r}=\frac{r^2-1}{2r(1+r)}=\frac12\Big(1-\frac1r\Big),
\]
and $1/r=\big(1+4e^{-(2-L)/X}\big)^{-1/2}$, as in Proposition~\ref{SM:prop:eos}. For the second statement,
group the words by $k$: $Z_\ell=\sum_k\Omega(\ell,k)e^{-(\ell L+(2-L)k+2)/X}$ has at most $\ell+1$ terms, so
$\frac1\ell\ln Z_\ell$ converges to the maximum over $x$ of $s(x)-v(x)/X$ (Laplace's method). Since $s$ is strictly concave, the
maximizer is unique and is characterized by $s'(x)=(2-L)/X$, which is again \eqref{SM:eq:compactivity}.
\end{proof}

At $X=\infty$ we have $a=b=1$, and $T$ becomes the Fibonacci matrix $\left(\begin{smallmatrix}1&1\\1&0\end{smallmatrix}\right)$ with
$\lambda=\varphi$. The identity $\ln\lambda(\infty)=\max s=\ln\varphi$ is Lecture~11's observation that ``the
maximal term carries the whole exponential growth''. Exercise~\ref{L11:ex:typical} asked for the
typical defect fraction $x^*$; we now see it as the $X=\infty$ point of a whole equation of state. Section~\ref{SM:sec:info}
identifies the measure behind it.

\subsection{Comparison with Bowles and Ashwin}

The same model was analysed independently, and in the same spirit, by Bowles and Ashwin (\emph{Edwards
entropy and compactivity in a model of granular matter}, Phys.\ Rev.~E 83, 2011). They use discs of
diameter $\sigma$ between walls a distance $H$ apart, with $H/\sigma<1+\sqrt{3/4}$, so that each disc
touches only its neighbours and one wall. With $\sigma=2$ and $H=w$ this is exactly our condition
$h<\sqrt3$. Their bond volumes $v_0=H\sqrt{(2-H)H}$ (cross-channel, $\sigma=1$) and $v_1=H$
(same wall) are our $L$ and $2$, rescaled by $\frac12$ and multiplied by the width. They count
$\Omega_J=\binom{N-M}{M}$ jammed states with $M$ defects among $N$ bonds, which differs from Lecture~11's
count only at the ends, and their entropy per particle and compactivity are exactly our $s(x)$ and
\eqref{SM:eq:compactivity}, in their units. They note that $X\to\infty$ at
$\theta=\frac12-\frac{\sqrt5}{10}$ and they restrict attention to $X>0$, remarking that denser-defect
packings would need negative compactivities. They also show that all the jammed packings are isostatic,
and that two subsystems exchanging defects equilibrate at equal compactivity. Ashwin and Bowles had
earlier (\emph{Complete jamming landscape of confined hard discs}, Phys.\ Rev.\ Lett.\ 102, 2009) treated the wider channel $1+\sqrt{3/4}<H/\sigma<2$ with a transfer
matrix on $32$ local tiles.

The open-ends caveat of Lecture~11 is precisely the boundary term on which the two counts differ; it
affects none of the limits above.

%======================================================================
\section{Information theory}
\label{SM:sec:info}

\subsection{Entropy and maximum entropy}

Lecture~7 counted coin flips and Lecture~11 maximized a multinomial coefficient. In both, the number of
sequences of length $N$ with letter frequencies $p$ is $e^{NH(p)+O(\ln N)}$, $H(p)=-\sum_ip_i\ln p_i$.
Shannon (1948) took $H$ as the measure of \emph{information}: such a source can be encoded with
$H(p)/\ln2$ bits per letter and no fewer. Jaynes (\emph{Information theory and statistical mechanics},
Phys.\ Rev.~106, 1957) turned the Lagrange multiplier computation of \S\ref{L11:sec:boltzmann} into a
principle: \emph{the least biased distribution compatible with what one knows has maximal entropy.}

\begin{proposition}[Maximum entropy gives Boltzmann]\label{SM:prop:jaynes}
Let $E_1,\dots,E_k$ be energies and $q_i=e^{-\beta E_i}/Z$ a Boltzmann distribution with mean energy
$\bar E=\sum q_iE_i$. Every probability vector $p$ with $\sum p_iE_i=\bar E$ satisfies $H(p)\le H(q)$, with
equality only for $p=q$.
\end{proposition}

\begin{proof}
The \emph{relative entropy} $D(p\Vert q)=\sum p_i\ln(p_i/q_i)$ is non-negative and vanishes only for $p=q$,
by Jensen's inequality for $\ln$. Since $-\ln q_i=\beta E_i+\ln Z$,
\[
  0\le D(p\Vert q)=-H(p)+\beta\sum_ip_iE_i+\ln Z=-H(p)+\beta\bar E+\ln Z=-H(p)+H(q).
\]
\end{proof}

Edwards' ansatz is the same principle with volume for energy: the flat measure on jammed states of
volume $V$ is the maximum entropy measure given $V$, and the canonical weight $e^{-V/X}$ is the maximum
entropy measure given the mean volume.

\subsection{The golden-mean channel}

Shannon's 1948 paper also treated \emph{constrained} channels, where not every sequence of symbols may be
sent. Its first example is the telegraph, whose symbols have different durations and may not follow each
other arbitrarily. The capacity of such a channel is the exponential growth rate of the number of
admissible messages. The simplest constraint is ``no two consecutive $1$'s'', the $(1,\infty)$
run-length-limited constraint of magnetic and optical recording. In symbolic dynamics the admissible
bi-infinite sequences form the \emph{golden-mean shift}. These are exactly the contact words of the
channel model.

\begin{figure}[ht]
\centering
\begin{tikzpicture}[scale=1.0,>={Stealth}]
  \node[circle,draw,thick,minimum size=0.8cm] (A) at (0,0) {$0$};
  \node[circle,draw,thick,minimum size=0.8cm] (B) at (2.6,0) {$1$};
  \draw[->,thick] (A) to[bend left=30] node[above,font=\small]{$1$} (B);
  \draw[->,thick] (B) to[bend left=30] node[below,font=\small]{$0$} (A);
  \draw[->,thick] (A) to[out=150,in=210,looseness=6] node[left,font=\small]{$0$} (A);
  \node[font=\small] at (6.2,0.3) {$T=\begin{pmatrix}1&1\\1&0\end{pmatrix}$,\quad $\lambda=\varphi$};
  \node[font=\small] at (6.2,-0.5) {Parry: $P(0\to1)=\varphi^{-2}$, $P(0\to0)=\varphi^{-1}$};
\end{tikzpicture}
\caption{The golden-mean shift. Admissible words are paths in this graph, where the state is the last
letter written. Its adjacency matrix is the Fibonacci matrix. With the weights $a=e^{-L/X}$, $b=e^{-2/X}$ on
the edges labelled $0$ and $1$, it becomes the Edwards transfer matrix of the channel.}
\label{SM:fig:golden}
\end{figure}

\begin{theorem}[Shannon 1948; Parry 1964]\label{SM:thm:golden}
\begin{enumerate}[nosep]
\item The number of binary words of length $n$ with no two consecutive $1$'s is $F_{n+2}$, and the
  capacity of the constraint is
  \[
    C=\lim_{n\to\infty}\frac1n\log_2F_{n+2}=\log_2\varphi\approx0.6942\ \text{bits per symbol}.
  \]
\item Among stationary Markov chains supported on admissible sequences, the entropy rate is maximized
  uniquely by the chain with $P(0\to1)=\varphi^{-2}$, $P(0\to0)=\varphi^{-1}$, $P(1\to0)=1$ (the
  \emph{Parry measure}). Its entropy rate is $\ln\varphi$ nats $=C$, and its frequency of $1$'s is
  $\frac{5-\sqrt5}{10}$.
\end{enumerate}
\end{theorem}

\begin{proof}
(1) The count is Lecture~11's Fibonacci argument. Equivalently it is $(1,1)T^{n-1}(1,1)^T$ with $T$ the
Fibonacci matrix, whose eigenvalues are $\varphi$ and $-\varphi^{-1}$, so the count is $c\varphi^n+O(\varphi^{-n})$.

(2) Such a chain is determined by $p=P(0\to1)$, since after a $1$ a $0$ is forced. Its stationary law
is $\pi_0=\frac1{1+p}$, $\pi_1=\frac p{1+p}$, and its entropy rate is $f(p)=\pi_0H(p)=H(p)/(1+p)$. Setting
$f'(p)=0$ gives $(1+p)\ln\frac{1-p}{p}=H(p)$, which simplifies to $2\ln(1-p)=\ln p$, that is
$(1-p)^2=p$, so $p=\frac{3-\sqrt5}2=\varphi^{-2}$ and $1-p=\varphi^{-1}$. Then
$H(p)=2p\ln\varphi+(1-p)\ln\varphi=(1+p)\ln\varphi$, so $f(p)=\ln\varphi$, and
$\pi_1=\frac{\varphi^{-2}}{1+\varphi^{-2}}=\frac1{\varphi^2+1}=\frac{5-\sqrt5}{10}$. Uniqueness: $f$ is
continuous on $[0,1]$, vanishes at $p=0$ and at $p=1$, and has only one critical point in $(0,1)$.
\end{proof}

(The variational principle of symbolic dynamics upgrades (2) to all shift-invariant measures: the Parry
measure is the unique measure of maximal entropy, and its entropy is the topological entropy
$\ln\varphi$. See Lind and Marcus in the Further reading.)

\begin{corollary}\label{SM:cor:capacity}
The maximal Edwards entropy per contact of the channel model equals the Shannon capacity of the
golden-mean constraint, $\ln\varphi$ nats $=\log_2\varphi$ bits. Edwards' flat measure on jammed
channel packings (the case $X=\infty$) is, in the limit, the Parry measure. More generally, the
canonical Edwards ensemble at compactivity $X$ is the Markov measure that maximizes entropy minus mean volume divided by $X$, for the weights $a,b$ of
Figure~\ref{SM:fig:golden}: an equilibrium (Gibbs) measure in the sense of dynamics (Exercise~\ref{SM:ex:markov}).
\end{corollary}

So the typical jammed state of Lecture~11, with its fraction $\frac{5-\sqrt5}{10}$ of same-wall contacts, is
the typical message of an optimally used golden-mean channel.

\begin{remark}[Capacity per unit length]
Shannon's telegraph asks for capacity per unit \emph{time} when symbols have durations $t_i$; for free
concatenation it is $\ln X_0$ with $\sum_iX_0^{-t_i}=1$. A no-$11$ word is a free concatenation of the blocks
$0$ and $10$, of channel lengths $L$ and $2+L$, so jammed configurations (any number of discs) of length
at most $t$ number $e^{C't+o(t)}$ with $e^{-C'L}+e^{-C'(L+2)}=1$. This is $\lambda(X)=1$ with $X=1/C'$, and by
Proposition~\ref{SM:prop:canon} $C'=\max_xs(x)/v(x)$. For $h=1.2$: $C'\approx0.2817$ at $x\approx0.266$, so
$X\approx3.55$.
\end{remark}

\subsection{Noise spheres: Shannon's packing argument}

Codes enter sphere packing through noise. A sender transmits a vector $x\in\R^n$ (one real number per
use of the channel) subject to a power constraint $\abs x^2\le nP$. The receiver sees $y=x+z$, where $z$ has
independent $N(0,N)$ coordinates. The capacity of this \emph{additive white Gaussian noise} channel is
the largest rate $R$ (bits per coordinate) such that $2^{nR}$ codewords can be chosen and decoded with
error probability tending to $0$ as $n\to\infty$.

\begin{theorem}[Shannon 1948, 1949]\label{SM:thm:awgn}
The capacity of the AWGN channel with signal power $P$ and noise power $N$ is
\[
  C=\tfrac12\log_2\Big(1+\frac PN\Big)\ \text{bits per coordinate}.
\]
\end{theorem}


\paragraph{Why packing?} The heuristic is the gas in a box of Lecture~7 again. By the law of large
numbers $\abs z^2=nN(1+O(n^{-1/2}))$: in high dimension the noise is concentrated on a thin shell of
radius $\sqrt{nN}$ around the transmitted codeword, just as the particle counts were concentrated near
their mean. It is also nearly orthogonal to $x$, so $\abs y^2\approx\abs x^2+\abs z^2\le n(P+N)$. Reliable
decoding requires the noise balls of radius $\sqrt{nN}$ around different codewords to be essentially
disjoint, and they all lie in a ball of radius about $\sqrt{n(P+N)}$. Comparing volumes gives at most
$\big(\sqrt{n(P+N)}/\sqrt{nN}\big)^n=(1+P/N)^{n/2}$ codewords, the bound in the theorem. Shannon's 1949 paper (\emph{Communication in the presence of noise}, Proc.\
IRE 37) turned this picture into a proof; achievability uses codewords chosen \emph{at random}, the
analogue of the non-constructive lower bounds of Lecture~1.

Demanding \emph{zero} error against every noise vector of length less than $\rho$ gives a genuine packing
problem. For codewords of equal power, on the sphere of radius $\sqrt{nP}$, this means angular separation
$\theta$ with $\sin(\theta/2)\ge\rho/\sqrt{nP}$, and the best such codes are counted by $M(n,\theta)$ of
Definition~\ref{L3:def:M}. So the Kabatiansky--Levenshtein bound of \S\ref{L3:sec:KL} is a coding theorem,
published as one (Problems of Information Transmission, 1978), and through Corollary~\ref{L3:cor:KL} it is
also the best known exponential upper bound on packing density.

\subsection{Perfect codes, Golay and Leech}

For binary codes (Definition~\ref{L04b:def:code}) the Hamming distance replaces the Euclidean one, and the
noise ball is a Hamming ball.

\begin{proposition}[Hamming's sphere packing bound]\label{SM:prop:hamming}
A binary code $C\subseteq\{0,1\}^n$ with minimal distance $d(C)\ge2e+1$ satisfies
\[
  \abs C\sum_{i=0}^e\binom ni\le2^n .
\]
\end{proposition}

\begin{proof}
The Hamming balls of radius $e$ around the codewords are disjoint, by the triangle inequality, and each
contains $\sum_{i\le e}\binom ni$ words.
\end{proof}

A code attaining equality is \emph{perfect}: its Hamming balls tile the cube, so every received word lies
within distance $e$ of exactly one codeword. The trivial perfect codes are a single word, the whole cube, and the
binary repetition codes of odd length. The nontrivial examples are rare and beautiful.
\begin{itemize}[nosep]
\item The \emph{Hamming codes} (1950): $n=2^r-1$, $e=1$, $\abs C=2^{n-r}$, and indeed
  $2^{n-r}(1+n)=2^n$.
\item The binary \emph{Golay code} (1949): $n=23$, $e=3$, $\abs C=2^{12}$, and
  \[
    2^{12}\Big(1+23+\binom{23}2+\binom{23}3\Big)=2^{12}(1+23+253+1771)=2^{12}\cdot2^{11}=2^{23}.
  \]
\item The ternary Golay code: $n=11$, $e=2$, $3^6$ words over $\mathbb F_3$, and
  $3^6(1+11\cdot2+\binom{11}2\cdot4)=3^6\cdot243=3^{11}$.
\end{itemize}
In fact these are all of them (Theorem~\ref{SM:thm:perfect} in the gallery).

Adding a parity bit turns the Hamming code of length $7$ into the extended Hamming code $[8,4,4]$, and
Construction~A (Definition~\ref{L04b:def:constrA}) turns that into $E_8$, as we saw in Lecture~6. Adding a
parity bit to the Golay code gives the extended Golay code $\mathcal G_{24}$, a $[24,12,8]$ code whose
$759$ words of weight $8$ (the \emph{octads}) form a Steiner system $S(5,8,24)$. Leech (1964, 1967) built
the Leech lattice from it. In the form given by Conway and Sloane (Ch.~4, \S11), $\sqrt8\,\Lambda_{24}$ is the set
of $x\in\Z^{24}$ such that, for some $m\in\{0,1\}$, all $x_i\equiv m\pmod2$, each set
$\{i:x_i\equiv\nu\pmod4\}$ ($\nu\in\Z/4$) is a word of $\mathcal G_{24}$, and $\sum_ix_i\equiv4m\pmod8$.
The minimal vectors have norm $32$ in these coordinates, and there are
\[
  \underbrace{\tbinom{24}2\cdot4}_{(\pm4,\pm4,0^{22})}+\underbrace{759\cdot2^7}_{(\pm2^8,0^{16})\text{ on an octad}}+\underbrace{24\cdot2^{12}}_{(\mp3,\pm1^{23})}
  =1104+97152+98304=196560
\]
of them, the kissing number in dimension $24$ quoted in Lecture~3. So $E_8$ and the Leech lattice, the two lattices now known to be optimal packings beyond dimension $3$,
come from the perfect Hamming code of length $7$ and the perfect Golay code of length $23$ by adding a
parity bit.

\subsection{Lattice codes}

Construction~A also runs the other way: a lattice $\Lambda\subset\R^n$ is a code for the Gaussian channel,
with infinitely many codewords. To respect the power constraint one keeps only the points of $\Lambda$ in a
bounded region. A good choice of region, which keeps the code a group, is the Voronoi cell of a coarser lattice $\Lambda'\subset\Lambda$,
which gives the \emph{nested lattice codes} $\Lambda/\Lambda'$. Erez and Zamir (IEEE Trans.\ Inform.\ Theory 50,
2004) proved that, with a random dither and a scaling at the receiver, nested lattice codes with
lattice decoding achieve the capacity $\frac12\log_2(1+P/N)$. Decoding is then
the nearest-point (Voronoi cell) problem of Lecture~1.

%======================================================================
\section{Gallery of theorems}
\label{SM:sec:gallery}

Results stated in the text (R\'enyi, Tonks and van Hove, Shannon, Kabatiansky--Levenshtein) are not
repeated here.

\begin{theorem}[Hales, Ann.\ of Math.\ 162, 2005; Viazovska, Ann.\ of Math.\ 185, 2017; Cohn--Kumar--Miller--Radchenko--Viazovska, Ann.\ of Math.\ 185, 2017]
The optimal sphere packing densities in dimensions $3$, $8$, $24$ are $\pi/\sqrt{18}$, $\pi^4/384$ and
$\pi^{12}/12!$. Equivalently, the hard-sphere gas of unit diameter in $\R^3$, $\R^8$, $\R^{24}$ has finite
free energy at number densities below, and infinite free energy above, $\sqrt2$, $16$ and $2^{24}$ respectively.
\end{theorem}
\noindent By Proposition~\ref{SM:prop:kepler}, $\rho_{\mathrm{cp}}$ is the optimal density divided by
$\vol\Ball^d(\frac12)$; it is the number density of FCC, $E_8$ and Leech, scaled to minimal distance~$1$.

\begin{theorem}[Hales--McLaughlin, J.~Amer.\ Math.\ Soc.\ 23, 2010]
In a packing of unit balls in $\R^3$, every Voronoi cell has volume at least that of a regular dodecahedron
of inradius $1$.
\end{theorem}
\noindent Fejes T\'oth's dodecahedral conjecture: the locally best density $\approx0.7547$ exceeds the global
optimum $0.7405$. This is geometric frustration as a theorem.

\begin{theorem}[Theil, Comm.\ Math.\ Phys.\ 262, 2006]
For a class of short-range, Lennard-Jones-like pair potentials in the plane, the ground-state energy per
particle of $N$ particles converges as $N\to\infty$ to the energy per particle of the triangular lattice:
pieces of the triangular lattice are asymptotically optimal.
\end{theorem}
\noindent Crystallization at zero temperature, the problem whose finite-temperature version for hard
spheres is open (\S\ref{SM:sec:hard}). Heitmann--Radin (1980) had done ``sticky discs''.

\begin{theorem}[Richthammer, Comm.\ Math.\ Phys.\ 345, 2016]
In the hard disc model in a $2n\times2n$ box with arbitrary boundary condition, the mean square
displacement of a particle near the centre is at least $c\log n$.
\end{theorem}
\noindent Hard discs have no long-range positional order: the most a two-dimensional ``crystal'' can
have is the hexatic orientational order seen numerically by Bernard and Krauth.

\begin{theorem}[Connelly--Gortler--Solomonides--Yampolskaya, Discrete Comput.\ Geom.\ 2020]
A collectively jammed packing of $n$ discs in a flat torus $\R^2/\Lambda$ whose radii and lattice are
generic has exactly $2n-1$ contacts, and its contact graph is isostatic.
\end{theorem}
\noindent The isostatic conjecture of Section~\ref{SM:sec:random} as a theorem in the plane (the lattice is
held fixed, and must itself be generic). Corollary~\ref{SM:cor:maxwell} gives the lower bound $2n-1$; the
content is the matching upper bound, $2n-2$ positions modulo translation plus the common scale of the
radii. The authors leave dimension $\ge3$ open.

\begin{theorem}[Cohn--Kumar, J.~Amer.\ Math.\ Soc.\ 20, 2007]
A configuration on $\Sph^{n-1}$ that has $m$ distinct distances and is a spherical $(2m-1)$-design
(``sharp'') minimizes the energy $\sum_{i\ne j}f(\abs{x_i-x_j}^2)$ among configurations of the same size,
for every completely monotonic $f$. The same holds for the $120$ vertices of the $600$-cell.
\end{theorem}
\noindent Examples: the icosahedron, the $240$ roots of $E_8$, the $196560$ minimal vectors of the Leech
lattice. These are zero-temperature ground states for a whole family of interactions at once, including
Frank's ideal $\{3,3,5\}$ order.

\begin{theorem}[Cohn--Kumar--Miller--Radchenko--Viazovska, Ann.\ of Math.\ 196, 2022]
The $E_8$ and Leech lattices are universally optimal among point configurations in $\R^8$ and $\R^{24}$:
they minimize energy for every potential that is a completely monotonic function of squared distance,
for example Gaussians $e^{-\alpha r^2}$ and inverse powers.
\end{theorem}
\noindent So the zero-temperature Gaussian core model is solved in dimensions $8$ and $24$ at every
density. The sphere packing theorems are recovered in the limit $\alpha\to\infty$.

\begin{theorem}[Tiet\"av\"ainen, SIAM J.~Appl.\ Math.\ 24, 1973; van Lint 1971; Zinov'ev--Leont'ev 1973]\label{SM:thm:perfect}
Every nontrivial perfect code over a finite field has the same length, number of words and minimal
distance as a Hamming code or one of the two Golay codes.
\end{theorem}
\noindent Tilings of the Hamming cube by Hamming balls are classified by their parameters (codes with
the Hamming parameters need not be linear). The Golay code's tiling is the
discrete shadow of the Leech lattice.

\begin{theorem}[Erez--Zamir, IEEE Trans.\ Inform.\ Theory 50, 2004]
Nested lattice codes with dithering, MMSE scaling and lattice decoding achieve the capacity
$\frac12\log_2(1+P/N)$ of the Gaussian channel.
\end{theorem}
\noindent Structured codes, decoded by finding the nearest lattice point, reach Shannon's capacity.

%======================================================================
\section*{Exercises}
\addcontentsline{toc}{section}{Exercises}


\begin{exercise}
(a) Solve \eqref{SM:eq:renyi} on $[2,3]$: show $M(x)=1+\frac{2(x-2)}{x-1}$ there. (b) Show that every saturated
packing of unit intervals in $[0,x]$ has at least $(x-1)/2$ intervals, and deduce $m\ge\frac12$ without
computing anything.
\end{exercise}

\begin{exercise}
(a) Show that the honeycomb disc packing is locally jammed but not collectively jammed: find a collective
motion. (b) Verify the claim of Figure~\ref{SM:fig:maxwell}.
\end{exercise}


\begin{exercise}\label{SM:ex:markov}
Let $r=(r_0,r_1)^T$ be a positive right eigenvector of $T=\left(\begin{smallmatrix}a&b\\a&0\end{smallmatrix}\right)$
for $\lambda=\lambda(X)$. Show that $P(i\to j)=T_{ij}r_j/(\lambda r_i)$ is a stochastic matrix, that its
stationary chain has frequency of $1$'s equal to $x(X)$, and that for $X=\infty$ it is the Parry measure
of Theorem~\ref{SM:thm:golden}.
\end{exercise}

\begin{exercise}
Compute the variance of the number of defects $k$ in the canonical ensemble to leading order in $\ell$, by
differentiating $\ln\lambda$ twice in $\ln b$. Show that it is $\ell\big/\!\big(-s''(x)\big)$, $x=x(X)$ (compare
Lemma~\ref{L11:lem:concave}).
\end{exercise}

\begin{exercise}
(a) Check that $1+90+\binom{90}2=2^{12}$, so the sphere packing bound allows a perfect binary code of
length $90$ with $2^{78}$ words and $e=2$. (b) Show that no such code exists. \emph{Hint:} assume
$0\in C$, so every word of weight $3$ lies at distance $\le2$ from a unique codeword, necessarily of
weight $5$. Count the weight-$3$ words that have $1$'s in the first two positions and are covered by
codewords of weight $5$ with $1$'s in the first two positions: each such codeword covers exactly three of them.
\end{exercise}

\begin{exercise}
(a) Let $z\in\R^n$ have independent $N(0,N)$ coordinates. Show that
$\abs z^2/n\to N$ in probability, and compare with the gas in a box of
\S\ref{L07:subsec:gas}. (b) Show that a fixed $x$ with $\abs x^2=nP$ satisfies $\ip xz/n\to0$ in
probability, and deduce $\abs{x+z}^2/n\to P+N$. (c) If disjoint balls of radius $\sqrt{n(N-\varepsilon)}$
lie in a ball of radius $\sqrt{n(P+N+\varepsilon)}$, bound their number by $2^{n(C+\delta(\varepsilon))}$.
\end{exercise}


\section*{Further reading}
\addcontentsline{toc}{section}{Further reading}
\begin{itemize}[nosep]
\item D.~Ruelle, \emph{Statistical Mechanics: Rigorous Results}, Benjamin, 1969.
\item S.~Torquato and F.~H.~Stillinger, \emph{Jammed hard-particle packings: From Kepler to Bernal and
  beyond}, Rev.\ Mod.\ Phys.\ 82 (2010).
\item A.~Baule, F.~Morone, H.~J.~Herrmann and H.~A.~Makse, \emph{Edwards statistical mechanics for jammed
  granular matter}, Rev.\ Mod.\ Phys.\ 90 (2018).
\item T.~M.~Cover and J.~A.~Thomas, \emph{Elements of Information Theory}, Wiley.
\item D.~Lind and B.~Marcus, \emph{An Introduction to Symbolic Dynamics and Coding}, Cambridge, 1995.
\item J.~H.~Conway and N.~J.~A.~Sloane, \emph{Sphere Packings, Lattices and Groups}, Springer, Chs.~1--5.
\end{itemize}
